% Options for packages loaded elsewhere
\PassOptionsToPackage{unicode}{hyperref}
\PassOptionsToPackage{hyphens}{url}
\documentclass[
  11pt,
]{article}
\usepackage{xcolor}
\usepackage[left=3cm,right=3cm,top=2.5cm,bottom=2.5cm]{geometry}
\usepackage{amsmath,amssymb}
\setcounter{secnumdepth}{5}
\usepackage{iftex}
\ifPDFTeX
  \usepackage[T1]{fontenc}
  \usepackage[utf8]{inputenc}
  \usepackage{textcomp} % provide euro and other symbols
\else % if luatex or xetex
  \usepackage{unicode-math} % this also loads fontspec
  \defaultfontfeatures{Scale=MatchLowercase}
  \defaultfontfeatures[\rmfamily]{Ligatures=TeX,Scale=1}
\fi
\usepackage{lmodern}
\ifPDFTeX\else
  % xetex/luatex font selection
\fi
% Use upquote if available, for straight quotes in verbatim environments
\IfFileExists{upquote.sty}{\usepackage{upquote}}{}
\IfFileExists{microtype.sty}{% use microtype if available
  \usepackage[]{microtype}
  \UseMicrotypeSet[protrusion]{basicmath} % disable protrusion for tt fonts
}{}
\usepackage{setspace}
\makeatletter
\@ifundefined{KOMAClassName}{% if non-KOMA class
  \IfFileExists{parskip.sty}{%
    \usepackage{parskip}
  }{% else
    \setlength{\parindent}{0pt}
    \setlength{\parskip}{6pt plus 2pt minus 1pt}}
}{% if KOMA class
  \KOMAoptions{parskip=half}}
\makeatother
\usepackage{longtable,booktabs,array}
\usepackage{caption}
\captionsetup[table]{skip=6pt}
\newcounter{none} % for unnumbered tables
\usepackage{calc} % for calculating minipage widths
% Correct order of tables after \paragraph or \subparagraph
\usepackage{etoolbox}
\makeatletter
\patchcmd\longtable{\par}{\if@noskipsec\mbox{}\fi\par}{}{}
\makeatother
% Allow footnotes in longtable head/foot
\IfFileExists{footnotehyper.sty}{\usepackage{footnotehyper}}{\usepackage{footnote}}
\makesavenoteenv{longtable}
\usepackage{graphicx}
\makeatletter
\newsavebox\pandoc@box
\newcommand*\pandocbounded[1]{% scales image to fit in text height/width
  \sbox\pandoc@box{#1}%
  \Gscale@div\@tempa{\textheight}{\dimexpr\ht\pandoc@box+\dp\pandoc@box\relax}%
  \Gscale@div\@tempb{\linewidth}{\wd\pandoc@box}%
  \ifdim\@tempb\p@<\@tempa\p@\let\@tempa\@tempb\fi% select the smaller of both
  \ifdim\@tempa\p@<\p@\scalebox{\@tempa}{\usebox\pandoc@box}%
  \else\usebox{\pandoc@box}%
  \fi%
}
% Set default figure placement to htbp
\def\fps@figure{htbp}
\makeatother
\setlength{\emergencystretch}{3em} % prevent overfull lines
\providecommand{\tightlist}{%
  \setlength{\itemsep}{0pt}\setlength{\parskip}{0pt}}
%% tools/stamp.tex
%% Provenance footer for PDFs rendered from this project.
%% Vendored by zzcollab; this file is not project-specific. The
%% footer prints the source document and the time of rendering.
%% The git version is defined here too, and so is available to a
%% project that wants it on the page, but it is left off the
%% default footer: it already names the staged copy in share/ and
%% appears in its manifest row. All values are supplied at render
%% time by tools/stamp-render.R, which writes a small generated
%% file that redefines the macros below. The \providecommand
%% defaults keep a bare LaTeX compile working when the document is
%% built without the wrapper.
\usepackage{fancyhdr}
\usepackage{xcolor}
\providecommand{\stampsource}{\detokenize{(source not recorded)}}
\providecommand{\stamptime}{(time not recorded)}
\providecommand{\stampversion}{\detokenize{(version not recorded)}}
\setlength{\footskip}{40pt}
\newcommand{\stampfootcontent}{%
  \footnotesize\color{gray}%
  \parbox[b]{\textwidth}{%
    \stamptime\hfill\thepage\\[2pt]%
    \ttfamily\raggedright\stampsource}}
\pagestyle{fancy}
\fancyhf{}
\fancyfoot[C]{\stampfootcontent}
\renewcommand{\headrulewidth}{0pt}
\renewcommand{\footrulewidth}{0.4pt}
%% The title page uses the `plain` page style; redefine it so the
%% stamp appears there too.
\fancypagestyle{plain}{%
  \fancyhf{}%
  \fancyfoot[C]{\stampfootcontent}%
  \renewcommand{\headrulewidth}{0pt}%
  \renewcommand{\footrulewidth}{0.4pt}}
\renewcommand{\stampsource}{\textasciitilde{}/\allowbreak{}prj/\allowbreak{}res/\allowbreak{}36-pmsimstats-ng/\allowbreak{}pmsimstats-ng/\allowbreak{}analysis/\allowbreak{}report/\allowbreak{}05-nof1-design-sensitivity/\allowbreak{}bullets.Rmd}
\renewcommand{\stamptime}{2026-08-15 16:49 PDT}
\renewcommand{\stampversion}{\detokenize{bc0b671-wip}}
\usepackage{amsmath}
\usepackage{amssymb}
\usepackage{booktabs}
\usepackage{longtable}
\usepackage{float}
\usepackage[mathlines]{lineno}
\linenumbers
\usepackage{bookmark}
\IfFileExists{xurl.sty}{\usepackage{xurl}}{} % add URL line breaks if available
\urlstyle{same}
% fallback for those not using the hyperref driver hyperxmp:
\makeatletter
\@ifundefined{xmpquote}{\newcommand{\xmpquote}[1]{#1}}{}
\makeatother
\hypersetup{
  pdftitle={Structured summary: Power sensitivity of a hybrid open-label-plus-blinded- discontinuation N-of-1 design under carryover{,} serial correlation{,} and variance-component variation},
  pdfauthor={pmsimstats team},
  hidelinks,
  pdfcreator={LaTeX via pandoc}}

\title{Structured summary: Power sensitivity of a hybrid open-label-plus-blinded- discontinuation N-of-1 design under carryover, serial correlation, and variance-component variation}
\author{pmsimstats team}
\date{August 15, 2026}

\begin{document}
\maketitle

{
\setcounter{tocdepth}{2}
\tableofcontents
}
\setstretch{1.1}
\emph{This document is a structured, section-by-section summary of the key points argued in the companion manuscript, ``Power sensitivity of a hybrid open-label-plus-blinded- discontinuation N-of-1 design under carryover, serial correlation, and variance-component variation.'' It is provided as a supplementary reading aid and does not stand on its own; all notation, derivations, simulation detail, and citations are in the main manuscript.}

\bigskip\noindent

\rule{\textwidth}{0.4pt}\bigskip

\section{Introduction}\label{introduction}

\textbf{Aggregated N-of-1 designs offer power advantages over parallel RCTs, but the conditions governing that advantage are incompletely characterized.}

\begin{itemize}\setlength{\itemsep}{2pt}\setlength{\parskip}{0pt}
\item N-of-1 designs accumulate within-patient evidence across repeated crossover periods, enabling meta-analytic population-level inference.
\item Recent simulations report N-of-1 procedures matching or exceeding RCT power at smaller sample sizes.
\item The parameter axes that govern when the N-of-1 advantage holds or reverses are not yet well specified.
\end{itemize}

\textbf{This paper reports a systematic sensitivity study of the Hendrickson hybrid N-of-1 design for the main drug effect estimand.}

\begin{itemize}\setlength{\itemsep}{2pt}\setlength{\parskip}{0pt}
\item The hybrid combines open-label titration, blinded discontinuation, and brief crossover phases.
\item Three comparators are evaluated: alternating A/B/A/B N-of-1, two-period crossover, and parallel-group RCT.
\item The target estimand is the main drug effect, distinct from the biomarker-treatment interaction target examined elsewhere.
\end{itemize}

\textbf{Twelve sensitivity sweeps address the gap left by single-axis simulation comparisons in the published literature.}

\begin{itemize}\setlength{\itemsep}{2pt}\setlength{\parskip}{0pt}
\item Published comparisons typically vary one or two parameters at a time, leaving the primary power-ordering axes unidentified.
\item The twelve sweeps span effect size, carryover half-life, decay-form misspecification, heterogeneity, within-patient variance, cycles, sample size, AR(1) correlation, period length, biomarker architecture, a misspecification factorial, and null calibration.
\item Each sweep axis is anchored to a specific precedent in the simulation-comparison literature.
\end{itemize}

\textbf{This paper extends the companion clinical comparison across twelve parameter axes for the hybrid design family.}

\begin{itemize}\setlength{\itemsep}{2pt}\setlength{\parskip}{0pt}
\item The companion paper provides a head-to-head comparison at the prazosin-calibrated baseline.
\item The present paper documents how that baseline comparison generalizes under systematic parameter perturbation.
\end{itemize}

\subsection{Aims}\label{aims}

\textbf{The study follows the ADEMP framework with three aims: power characterization, design-ordering identification, and Type I error verification.}

\begin{itemize}\setlength{\itemsep}{2pt}\setlength{\parskip}{0pt}
\item Aim 1: characterize the hybrid's power profile under systematic perturbation of published comparison axes.
\item Aim 2: identify parameter regimes where the hybrid is preferred over the parallel RCT and where the ordering reverses.
\item Aim 3: verify Type I error control with MCSE below 1\%.
\end{itemize}

\subsection{Pipeline}\label{pipeline}

\textbf{All sweeps use a baseline of $N = 35$ participants and 8 weekly timepoints, with sample size matched per design rather than per observation count.}

\begin{itemize}\setlength{\itemsep}{2pt}\setlength{\parskip}{0pt}
\item The pipeline is implemented in the \texttt{nof1power} R package, which re-implements the Hendrickson 2020 reference pipeline.
\item The $N = 35$ minimum and per-design matching convention follow Blackston 2019 and Hendrickson 2020.
\item Parameters are varied from this baseline one axis at a time across the twelve sweeps.
\end{itemize}

\subsection{Dual data-generating architectures}\label{dual-data-generating-architectures}

\textbf{Two data-generating architectures encode the biomarker interaction differently, producing substantively different power profiles under carryover.}

\begin{itemize}\setlength{\itemsep}{2pt}\setlength{\parskip}{0pt}
\item The covariance-moderation architecture (MVN) encodes the interaction in the covariance structure; on-drug correlation decays exponentially off-drug.
\item The mean-moderation architecture shifts bio-response means additively at on-drug timepoints by the standardized biomarker.
\item The two architectures are crossed in sweep S10; the remaining eleven sweeps use the covariance-moderation architecture as the default.
\end{itemize}

\subsection{Response curves and carryover}\label{response-curves-and-carryover}

\textbf{Latent response factors follow a modified Gompertz curve; carryover is applied via a recursive decay from an extensible decay-family.}

\begin{itemize}\setlength{\itemsep}{2pt}\setlength{\parskip}{0pt}
\item The Gompertz curve passes through the origin and asymptotes to \texttt{maxr}, parameterized by displacement and rate.
\item Off-drug carryover is applied recursively using one of four decay families: exponential, linear, Weibull, or power.
\item The DGP decay family and the analysis-side assumption can differ; sweep S3 quantifies the cost of that mismatch.
\end{itemize}

\subsection{Covariance structure and analysis}\label{covariance-structure-and-analysis}

\textbf{The analysis model is an \texttt{nlme::lme} with \texttt{corCAR1} residual correlation; the target estimand is the \texttt{Dbc} main drug effect coefficient.}

\begin{itemize}\setlength{\itemsep}{2pt}\setlength{\parskip}{0pt}
\item The covariance matrix couples within-factor AR(1), cross-factor correlations, and biomarker differential correlation; non-PD matrices are corrected with a per-cell flag.
\item The fixed-effects formula is \texttt{Sx \textasciitilde{} bm + t + Dbc + bm:Dbc} with a random intercept per participant.
\item Singular fits are detected from the random-effects variance and flagged per replicate.
\end{itemize}

\subsection{Reproducibility and execution}\label{reproducibility-and-execution}

\textbf{Per-replicate rather than pre-aggregated output is written, enabling Morris MCSE computation of all performance measures at render time.}

\begin{itemize}\setlength{\itemsep}{2pt}\setlength{\parskip}{0pt}
\item Each sweep seeds via \texttt{set.seed(42 + sweep\_id)} and uses furrr multisession parallelism.
\item RDS artifacts are written to \texttt{analysis/data/derived\_data/sensitivity/} and orchestrated by \texttt{run-all.R}.
\item Per-replicate storage allows power, bias, MSE, empirical SE, and Type I error to be computed with MCSE at render time.
\end{itemize}

\section{Results}\label{results}

\textbf{All twelve sweeps completed in approximately 32 minutes; per-cell metrics are computed at render time from per-replicate RDS artifacts.}

\begin{itemize}\setlength{\itemsep}{2pt}\setlength{\parskip}{0pt}
\item Wall-clock time was approximately 32 minutes on 8 cores with furrr multisession.
\item Summary metrics are derived at render time from \texttt{analysis/data/derived\_data/sensitivity/}.
\item The target estimand throughout is the \texttt{Dbc} main drug effect coefficient.
\end{itemize}

\textbf{The hybrid is the primary design; the alternating A/B/A/B variant is retained as a sensitivity comparator for sweeps that vary cycle count or period length.}

\begin{itemize}\setlength{\itemsep}{2pt}\setlength{\parskip}{0pt}
\item The hybrid's fixed eight-timepoint structure does not scale with cycle count or period length.
\item The alternating variant is used in S6 and S9 where that scalability is required.
\end{itemize}

\textbf{Across sweeps, the hybrid outperforms the parallel RCT on sample size, variance robustness, and serial-correlation tolerance; the RCT is more robust to long carryover.}

\begin{itemize}\setlength{\itemsep}{2pt}\setlength{\parskip}{0pt}
\item The hybrid reaches 80\% power at smaller sample sizes, under higher within-patient variance, and under stronger AR(1) than the parallel RCT.
\item The parallel RCT is more robust to carryover exceeding the within-subject period length and to decay-form misspecification.
\item Type I error is near nominal for the hybrid and alternating variant; the two-period crossover shows robust inflation of 0.10 to 0.13.
\end{itemize}

\subsection{S1. Effect-size sweep}\label{s1.-effect-size-sweep}

\textbf{The hybrid reaches 80\% power at $\delta \approx -1$, a half-unit smaller effect size than the parallel RCT requires.}

\begin{itemize}\setlength{\itemsep}{2pt}\setlength{\parskip}{0pt}
\item Hybrid power rises from 0.066 at the null to saturation at 1.000 for $|\delta| \geq 2$.
\item Parallel RCT power follows the same shape at a lower level, reaching 0.941 at $\delta = -2$.
\item The hybrid's 80\%-power effect-size threshold is approximately 0.5 units smaller than the RCT's.
\end{itemize}

\textbf{Type I error is marginally above nominal for the hybrid and marginally below for the RCT; neither is grossly miscalibrated, but the hybrid's upward departure warrants attention.}

\begin{itemize}\setlength{\itemsep}{2pt}\setlength{\parskip}{0pt}
\item Hybrid Type I is 0.066 (95\% CI [0.050, 0.082]); parallel RCT Type I is 0.034 (95\% CI [0.022, 0.046]).
\item Both confidence intervals cover 0.05 at the margin.
\item The hybrid's modest upward departure merits caution when the study is powered near the 5\% boundary.
\end{itemize}

\subsection{S2. Carryover half-life sweep}\label{s2.-carryover-half-life-sweep}

\textbf{The hybrid tolerates carryover up to 5 days before power erodes; the alternating variant degrades from 0.90 to 0.30 across the half-life range.}

\begin{itemize}\setlength{\itemsep}{2pt}\setlength{\parskip}{0pt}
\item Hybrid power is 1.000 through $t_{1/2} = 5$ days, declining to 0.736 at 14 days.
\item The alternating N-of-1 is more sensitive, dropping from 0.896 to 0.304 across the same range.
\item The parallel RCT is essentially immune, with power flat at 0.92 to 0.96 across the grid.
\end{itemize}

\textbf{The hybrid's carryover robustness is attributable to its open-label titration phase; the RCT ordering reversal begins when $t_{1/2}$ exceeds the period length.}

\begin{itemize}\setlength{\itemsep}{2pt}\setlength{\parskip}{0pt}
\item The four-week open-label titration accumulates treatment signal that is not erased by the post-discontinuation off-drug period.
\item The alternating design relies on repeated within-subject transitions and sheds signal proportional to the unwashed residual.
\item The hybrid-versus-RCT power ordering crossover begins at approximately $t_{1/2} = 10$ days, consistent with Blackston 2019.
\end{itemize}

\subsection{S3. Decay-form misspecification sweep}\label{s3.-decay-form-misspecification-sweep}

\textbf{The decay-form sweep uses the alternating variant because the hybrid's multi-phase \texttt{tsd} pattern produces non-PD matrices under the power decay form.}

\begin{itemize}\setlength{\itemsep}{2pt}\setlength{\parskip}{0pt}
\item Five DGP decay forms are tested while the analysis always assumes exponential decay.
\item The hybrid is excluded because its multi-phase \texttt{tsd} structure interacts poorly with the power form.
\item The alternating variant provides a clean test of decay-form misspecification sensitivity.
\end{itemize}

\textbf{Main-effect power ranges only 15 percentage points across decay forms, substantially less sensitive to misspecification than the interaction estimand.}

\begin{itemize}\setlength{\itemsep}{2pt}\setlength{\parskip}{0pt}
\item Power ranges from 0.754 (Weibull $k=0.8$) to 0.904 (linear) under the alternating N-of-1.
\item The 15-point range is smaller than under the interaction target, indicating less sensitivity to decay-form misspecification.
\item The parallel RCT is flat at 0.93 to 0.95 across all forms, consistent with Gärtner 2023.
\end{itemize}

\subsection{S4. Between-patient heterogeneity sweep}\label{s4.-between-patient-heterogeneity-sweep}

\textbf{The hybrid is saturated across the $\tau$ range; the parallel RCT loses 12 percentage points as between-patient heterogeneity doubles.}

\begin{itemize}\setlength{\itemsep}{2pt}\setlength{\parskip}{0pt}
\item Hybrid power is 1.000 at all six $\tau$ levels with stable mean $\beta \approx -0.94$.
\item Parallel RCT power declines from 0.998 at $\tau = 0.5$ to 0.874 at $\tau = 2.0$.
\item The RCT loses approximately 12 percentage points as $\tau$ doubles from 1.0 to 2.0.
\end{itemize}

\textbf{The Senn 2018 variance-decomposition prediction is reproduced: $\tau^2$ enters the RCT estimator but is absorbed by the within-subject contrast in N-of-1 designs.}

\begin{itemize}\setlength{\itemsep}{2pt}\setlength{\parskip}{0pt}
\item Between-patient heterogeneity contributes to $\mathrm{Var}(\hat\Delta)$ in the parallel RCT but not in the within-subject N-of-1 contrast.
\item The hybrid's saturation reflects the prazosin-calibrated baseline where effect size is large relative to within-patient variance.
\item At smaller effect sizes or larger $\sigma$, the N-of-1 advantage would widen rather than flatten as $\tau$ increases.
\end{itemize}

\subsection{S5. Within-patient variance sweep}\label{s5.-within-patient-variance-sweep}

\textbf{The hybrid is essentially saturated across the within-patient variance range; the parallel RCT loses approximately 18 percentage points.}

\begin{itemize}\setlength{\itemsep}{2pt}\setlength{\parskip}{0pt}
\item Hybrid power declines from 1.000 to 0.986 across $\sigma \in \{0.5, \ldots, 1.5\}$, a one-percentage-point loss.
\item Parallel RCT power declines from 0.996 to 0.814 across the same range, an 18-point loss.
\item The ordering is consistent with the N-of-1 design summing $k$ paired contrasts per patient.
\end{itemize}

\textbf{The hybrid's quasi-saturation is specific to the prazosin-calibrated regime; at smaller effect sizes the $\sigma$ dependence would be steeper.}

\begin{itemize}\setlength{\itemsep}{2pt}\setlength{\parskip}{0pt}
\item The eight-timepoint structure plus open-label titration provides sufficient signal accumulation to absorb substantial within-patient noise.
\item The Senn 2018 prediction is reproduced: N-of-1 absorbs within-patient variance via $k$ paired contrasts.
\item Saturation should not be extrapolated beyond the prazosin-calibrated effect-size regime.
\end{itemize}

\subsection{S6. Cycles-per-patient sweep}\label{s6.-cycles-per-patient-sweep}

\textbf{Alternating N-of-1 power rises concavely with $k$, from 0.294 at $k = 2$ to 0.980 at $k = 16$, as predicted by the Senn 2018 decomposition.}

\begin{itemize}\setlength{\itemsep}{2pt}\setlength{\parskip}{0pt}
\item The alternating variant is used because the hybrid's fixed structure does not scale with $k$.
\item Power rises from 0.294 at $k = 2$ to 0.804 at $k = 8$ (baseline) to 0.980 at $k = 16$.
\item The concave shape is consistent with the Senn 2018 variance-decomposition prediction.
\end{itemize}

\textbf{For the main-effect target, increasing $k$ slightly favors the RCT; the hybrid's advantage over the alternating variant comes from the open-label phase, not cycle count.}

\begin{itemize}\setlength{\itemsep}{2pt}\setlength{\parskip}{0pt}
\item At $k = 2$ the alternating N-of-1 and the RCT are in a dead heat; at $k \geq 6$ the RCT surpasses the alternating N-of-1 by up to 14 points.
\item Each additional timepoint contributes a between-subject observation for the RCT but averages into already-accumulated paired contrasts for the alternating N-of-1.
\item The hybrid's primary power advantage over the alternating design arises from the open-label titration phase, not from the number of cycles.
\end{itemize}

\subsection{S7. Patient-count sweep}\label{s7.-patient-count-sweep}

\textbf{The hybrid saturates at high power by $N = 16$; the parallel RCT reaches the same threshold between $N = 24$ and $N = 35$.}

\begin{itemize}\setlength{\itemsep}{2pt}\setlength{\parskip}{0pt}
\item Hybrid power is high at modest sample sizes and saturates rapidly with $N$.
\item The parallel RCT follows a similar but offset trajectory, reaching 80\% power between $N = 24$ and $N = 35$.
\item The grid anchors the project-standard $N = 35$ minimum for cross-paper comparison.
\end{itemize}

\textbf{The hybrid achieves an equivalent power target at approximately one-third to one-half the parallel-RCT sample size, consistent with Blackston 2019 and Carrozzo 2025.}

\begin{itemize}\setlength{\itemsep}{2pt}\setlength{\parskip}{0pt}
\item The hybrid reaches 80\% power at $N \leq 16$; the parallel RCT reaches it between $N = 24$ and $N = 35$.
\item The sample-size ratio of one-third to one-half matches the Blackston 2019 N-of-1 efficiency advantage.
\item This is consistent with Carrozzo 2025, in which the meta-analytic N-of-1 required fewer participants across most heterogeneity and carryover settings.
\end{itemize}

\subsection{S8. AR(1) serial-correlation sweep}\label{s8.-ar1-serial-correlation-sweep}

\textbf{The hybrid is insensitive to AR(1) correlation; the parallel RCT loses 24 percentage points of power as $\rho$ rises from 0 to 0.9.}

\begin{itemize}\setlength{\itemsep}{2pt}\setlength{\parskip}{0pt}
\item Hybrid power is 1.000 across $\rho \in \{0, 0.3, 0.5, 0.7, 0.9\}$.
\item Parallel RCT power declines from 0.984 to 0.746, and mean $\hat\beta$ shrinks from $-1.01$ to $-0.57$.
\item The RCT loss is monotonic with $\rho$ across the examined range.
\end{itemize}

\textbf{The AR(1) asymmetry arises because the \texttt{corCAR1} structure absorbs correlation within subjects but cannot manufacture the missing within-subject contrasts for the parallel RCT.}

\begin{itemize}\setlength{\itemsep}{2pt}\setlength{\parskip}{0pt}
\item Strong AR(1) produces coherent within-subject trajectories for the hybrid, absorbed by \texttt{corCAR1} with negligible power loss.
\item For the RCT, strong AR(1) inflates between-subject mean contrast variance via effective-sample-size shrinkage.
\item This is consistent with Wang and Schork 2019: N-of-1 designs are insensitive to serial correlation, while parallel-group designs are not.
\end{itemize}

\subsection{S9. Period-length sweep}\label{s9.-period-length-sweep}

\textbf{Main-effect power increases monotonically with period length, from 0.40 at 3 days to 1.00 at 14 days; carryover half-life has only a modest effect at fixed period length.}

\begin{itemize}\setlength{\itemsep}{2pt}\setlength{\parskip}{0pt}
\item Power rises from 0.40 to 0.44 at 3-day periods to 0.997 to 1.00 at 14-day periods.
\item At fixed period length, varying carryover half-life across 0.5 to 7 days produces only modest power differences.
\item The sweep uses the alternating variant, as the hybrid's fixed structure does not scale with period length.
\end{itemize}

\textbf{For the main-effect target, longer periods are preferred, reversing the interaction-target recommendation; carryover should not exceed approximately 50\% of the period length.}

\begin{itemize}\setlength{\itemsep}{2pt}\setlength{\parskip}{0pt}
\item Longer periods accumulate more Gompertz drug-response signal per timepoint.
\item The practical guideline is the reverse of the interaction-target recommendation: longer periods benefit main-effect detection.
\item Wang and Schork 2019 note this tradeoff explicitly: frequent alternation benefits correlation-modeled analyses but not absolute-effect inference.
\end{itemize}

\subsection{S10. Biomarker interaction strength crossed with DGP architecture}\label{s10.-biomarker-interaction-strength-crossed-with-dgp-architecture}

\textbf{Main-effect power is invariant to $c_{bm}$ and DGP architecture, confirming that the main drug effect is determined by Gompertz accumulation, not biomarker coupling.}

\begin{itemize}\setlength{\itemsep}{2pt}\setlength{\parskip}{0pt}
\item Hybrid power is 1.000 across all 16 $c_{bm}$-by-architecture-by-halflife combinations.
\item Alternating N-of-1 power varies from 0.682 to 0.916; parallel RCT power varies from 0.914 to 0.958.
\item The flat main-effect power across $c_{bm}$ values provides a robustness check on the DGP implementation.
\end{itemize}

\textbf{The flat $c_{bm}$ power curves are a DGP implementation check: when the main effect is the target, biomarker coupling should not affect power, and it does not.}

\begin{itemize}\setlength{\itemsep}{2pt}\setlength{\parskip}{0pt}
\item The main drug effect is determined by Gompertz $\mathrm{max}[\mathrm{br}]$ and on-drug accumulation, not by biomarker coupling.
\item Approximately flat power across $c_{bm}$ for all three designs confirms the DGP correctly separates the two signal sources.
\item The reader should not conflate this main-effect invariance with the biomarker-interaction target, where $c_{bm}$ is the primary power driver.
\end{itemize}

\subsection{S11. Misspecification factorial}\label{s11.-misspecification-factorial}

\textbf{The parallel RCT shows statistically robust Type I inflation of 0.112 to 0.126 under high DGP AR(1); the hybrid is controlled near nominal across all four cells.}

\begin{itemize}\setlength{\itemsep}{2pt}\setlength{\parskip}{0pt}
\item Hybrid null Type I is 0.036 to 0.056, within 2 MCSE of nominal across all four DGP combinations.
\item Parallel RCT Type I is 0.018 to 0.024 (conservative) when no AR(1) is present, and 0.112 to 0.126 (inflated) under $\rho = 0.7$.
\item The inflation is statistically robust at MCSE 0.014 to 0.015.
\end{itemize}

\textbf{The RCT Type I inflation arises because the between-subject variance estimator understates the true standard error when DGP serial correlation is strong.}

\begin{itemize}\setlength{\itemsep}{2pt}\setlength{\parskip}{0pt}
\item Under strong DGP AR(1), \texttt{corCAR1} cannot manufacture the missing within-subject contrasts for a between-subject comparison.
\item The hybrid is immune because its primary treatment contrast is within-subject.
\item The mechanism is analogous to the Blackston 2019 carryover Type I inflation finding, though the source here is AR(1) rather than carryover.
\end{itemize}

\subsection{S12. Null-effect calibration}\label{s12.-null-effect-calibration}

\textbf{Hybrid and alternating N-of-1 Type I error is near nominal; the two-period crossover shows robust inflation of 0.100 to 0.130 attributable to underidentified \texttt{corCAR1}.}

\begin{itemize}\setlength{\itemsep}{2pt}\setlength{\parskip}{0pt}
\item Hybrid Type I: [0.056, 0.062]; alternating N-of-1: [0.057, 0.064]; both within 2 MCSE of nominal.
\item Parallel RCT: [0.029, 0.060], three cells conservative and one nominal.
\item Two-period crossover: [0.100, 0.130], statistically robust inflation at MCSE 0.006 to 0.007.
\end{itemize}

\textbf{The two-period crossover should not be analyzed with \texttt{corCAR1}; at two within-subject timepoints the parameter is underidentified and the rejection rate inflates by a practically meaningful amount.}

\begin{itemize}\setlength{\itemsep}{2pt}\setlength{\parskip}{0pt}
\item Two within-subject timepoints cannot identify the \texttt{corCAR1} parameter reliably.
\item The inflation of 0.100 to 0.130 is statistically robust and practically meaningful.
\item Alternative analyses (paired $t$-test on change scores or lme with unstructured residual) are recommended for this design.
\end{itemize}

\subsection{Synthesis}\label{synthesis}

\textbf{Five substantive findings organize the twelve sweeps: sample-size advantage, variance robustness, carryover design-dependence, AR(1) asymmetry, and Type I error calibration.}

\begin{itemize}\setlength{\itemsep}{2pt}\setlength{\parskip}{0pt}
\item The hybrid reaches 80\% power at $N \leq 16$; the parallel RCT requires $N = 24$ to $35$, a one-third to one-half sample-size ratio.
\item The hybrid is saturated across $\sigma$ and $\tau$ sweeps; the RCT degrades by 18 and 12 percentage points respectively.
\item Carryover sensitivity is design-specific: the hybrid tolerates $t_{1/2} \leq 5$ days; the alternating variant drops from 0.90 to 0.30 across the range; the RCT is immune.
\end{itemize}

\subsubsection{Hendrickson et al.~(2020): predictive biomarker validation}\label{hendrickson-et-al.-2020-predictive-biomarker-validation}

\textbf{Hendrickson 2020 showed hybrid power collapsing from 0.95 to 0.18 at the biomarker-interaction target under a 0.7-day carryover half-life.}

\begin{itemize}\setlength{\itemsep}{2pt}\setlength{\parskip}{0pt}
\item At $N = 35$, $c_{bm} = 0.6$, no carryover, biomarker-interaction power was 0.95 (hybrid), 0.96 (crossover, CO), 0.94 (open-label plus blinded discontinuation, OL+BDC), and 0.25 (open-label, OL).
\item A half-life of only 0.7 days collapsed hybrid power to 0.18, demonstrating extreme carryover sensitivity at the interaction estimand.
\item This contrasts sharply with the present study's main-effect findings, where hybrid power is 1.000 through $t_{1/2} = 5$ days.
\end{itemize}

\textbf{The present study extends Hendrickson via Morris MCSE reporting, decay-form misspecification (S3), and main-effect carryover and period-length quantification (S2, S6, S9).}

\begin{itemize}\setlength{\itemsep}{2pt}\setlength{\parskip}{0pt}
\item Design ranking at the biomarker-interaction target reproduces Hendrickson by construction.
\item The key substantive divergence is estimand-specific: main-effect hybrid power is saturated through 5-day carryover, while interaction-target power collapses at under 1 day.
\item Both findings are correct; they describe different estimands and must not be conflated.
\end{itemize}

\subsubsection{Wang and Schork (2019): AR(1), heteroscedasticity, and period stacking}\label{wang-and-schork-2019-ar1-heteroscedasticity-and-period-stacking}

\textbf{Wang and Schork 2019 show roughly linear power erosion with $\rho$ for single-period designs; period stacking substantially mitigates this loss.}

\begin{itemize}\setlength{\itemsep}{2pt}\setlength{\parskip}{0pt}
\item For $1p \times 2i$, power erodes from 0.851 at $\rho = 0$ to 0.126 at $\rho = 0.75$.
\item For $40p \times 2i$, the same $\rho$ range produces only a 0.17-point power loss, confirming period stacking as a mitigation.
\end{itemize}

\textbf{S8 is the aggregated-level analogue of Wang and Schork: the hybrid absorbs AR(1) via \texttt{corCAR1}; the parallel-RCT power loss directionally matches Wang and Schork's between-period profile.}

\begin{itemize}\setlength{\itemsep}{2pt}\setlength{\parskip}{0pt}
\item Hybrid power is saturated at 1.00 across $\rho \in \{0, 0.3, 0.5, 0.7, 0.9\}$ at the aggregated level.
\item Parallel RCT power drops from 0.98 to 0.75, directionally consistent with Wang and Schork's between-period power loss.
\item The S11 Type I inflation of 0.11 to 0.13 is consistent with the Wang and Schork mechanism, though magnitudes depend on small-sample instability of \texttt{corCAR1} at $N = 35$.
\end{itemize}

\subsection{Limitations}\label{limitations}

\textbf{Four limitations apply: hybrid saturation at the prazosin baseline, alternating-variant substitution in three sweeps, analysis-only \texttt{corCAR1} in S8, and small-sample instability in the S11 RCT Type I estimate.}

\begin{itemize}\setlength{\itemsep}{2pt}\setlength{\parskip}{0pt}
\item Saturation in S2, S4, S5, S8, S10, and S11 prevents quantifying the hybrid's relative advantage in sub-ceiling regimes; sweeps at smaller effect sizes are the priority extension.
\item S3, S6, and S9 use the alternating variant because the hybrid cannot scale with decay form, cycle count, or period length; results apply to the alternating family, not directly to the hybrid.
\item The S11 RCT Type I ranges (0.112 to 0.126 vs. 0.029 to 0.060 in S12) are not contradictory; they reflect different DGP regimes.
\end{itemize}

\subsection{N-of-1 design variants not evaluated}\label{n-of-1-design-variants-not-evaluated}

\textbf{Several N-of-1 design variants are not evaluated; these fall into designs expected to be weaker than those reported here and designs expected to be stronger.}

\begin{itemize}\setlength{\itemsep}{2pt}\setlength{\parskip}{0pt}
\item Unevaluated designs include single-subject, OL-only, OL+BDC-only, full multi-cycle crossover, complete crossover, sequential-monitoring, and randomized-block variants.
\item Single-subject and reduced variants are expected to produce power at or below the reported numbers; results are an upper bound.
\item Full multi-cycle, complete-design, and sequential-monitoring variants are expected to exceed the reported numbers; results are a lower bound.
\end{itemize}

\textbf{The present results bound the unevaluated variants: weaker designs are upper-bounded and stronger designs are lower-bounded by the reported numbers.}

\begin{itemize}\setlength{\itemsep}{2pt}\setlength{\parskip}{0pt}
\item Simpler variants (single-subject, OL-only, OL+BDC-only, two-cycle crossover) are expected at or below the reported power.
\item More powerful variants (multi-cycle, complete-design, sequential-monitoring, randomized-block) are expected at or above the reported power.
\item Extending the present framework to cover these variants is a tractable implementation task rather than a conceptual redesign.
\end{itemize}

\subsection{Prototype Monte Carlo confirmation}\label{prototype-monte-carlo-confirmation}

\textbf{A 500-replicate prototype ran in 191 seconds with 100\% convergence; the MCSE on a power estimate near 0.3 is approximately 0.020.}

\begin{itemize}\setlength{\itemsep}{2pt}\setlength{\parskip}{0pt}
\item The 15-cell $c_{br} \times N$ prototype grid was completed in 191 seconds on 9 cores.
\item Convergence was 100\% across all 7,500 fits.
\item Per-replicate output and Morris ADEMP summary are at \texttt{analysis/data/quick-sim/}.
\end{itemize}

\textbf{The apparent non-monotonic power dependence on $c_{br}$ from the 100-replicate prototype does not survive at 500 replicates; the signal is broadly flat.}

\begin{itemize}\setlength{\itemsep}{2pt}\setlength{\parskip}{0pt}
\item At $N = 25$, the five-level swing is 0.024 (one MCSE), consistent with sampling noise around a constant.
\item At $N = 35$ and $N = 50$, the peak-to-trough swings of approximately three MCSE are suggestive but not decisive, and the peaks do not coincide on $c_{br}$.
\item The 100-replicate dip at $c_{br} = 0.15$ has disappeared, confirming it was sampling noise.
\end{itemize}

\textbf{The dominant axis is sample size, not $c_{br}$; no prototype cell attains 80\% power, and the grid does not bound the 80\%-power region.}

\begin{itemize}\setlength{\itemsep}{2pt}\setlength{\parskip}{0pt}
\item Mean power rises from 0.18 at $N = 25$ to 0.31 at $N = 50$; the $c_{br}$ effect is a possible shallow optimum near $[0.30, 0.45]$.
\item No $(c_{br}, N)$ cell in the prototype reaches 80\% power; a production run requires larger $N$, larger $br_{\max}$, or both.
\item The prototype grid does not bound the 80\%-power region, which is the central deliverable for trial planning.
\end{itemize}

\section{Conclusions}\label{conclusions}

\textbf{Three conclusions: the hybrid is a competent main-effect default, design ordering is estimand-dependent, and two-period crossovers should not be analyzed with \texttt{corCAR1}.}

\begin{itemize}\setlength{\itemsep}{2pt}\setlength{\parskip}{0pt}
\item The hybrid absorbs variance and serial correlation, tolerates carryover up to five days, and achieves 80\% power at smaller sample sizes than the parallel RCT.
\item Design ordering is substantially estimand-dependent: the hybrid's advantage is larger for biomarker interactions than for main effects.
\item The two-period crossover's robust Type I inflation of 0.10 to 0.13 warrants a change to standard analytic practice when \texttt{corCAR1} is fitted at two within-subject timepoints.
\end{itemize}

\textbf{The hybrid N-of-1 is supported as a default candidate for trial planning in chronic neurobehavioral conditions where moderate treatment heterogeneity and short-to-moderate carryover are expected.}

\begin{itemize}\setlength{\itemsep}{2pt}\setlength{\parskip}{0pt}
\item The companion clinical paper provides the baseline head-to-head comparison; this paper documents its generalization across twelve parameter axes.
\item Robustness across most axes, with the identified exceptions for long carryover and two-period crossover analysis, supports the hybrid as a default candidate.
\end{itemize}

\section{Data availability statement}\label{data-availability-statement}

\textbf{Per-replicate Monte Carlo outputs and ADEMP cell-level summaries are versioned in the pmsimstats-ng repository; exact paths are in the Reproducibility section.}

\begin{itemize}\setlength{\itemsep}{2pt}\setlength{\parskip}{0pt}
\item Monte Carlo outputs are under \texttt{analysis/data/}.
\item Simulation driver scripts are under \texttt{analysis/scripts/}.
\end{itemize}

\section{Reproducibility}\label{reproducibility}

\textbf{The manuscript was rendered with R 4.4.x inside the zzcollab Docker container; principal packages are \texttt{nlme}, \texttt{parallel}, \texttt{data.table}, \texttt{furrr}, and \texttt{purrr}.}

\begin{itemize}\setlength{\itemsep}{2pt}\setlength{\parskip}{0pt}
\item Image hash is in \texttt{Dockerfile}; package set is in \texttt{renv.lock}.
\item \texttt{nlme} provides \texttt{corCAR1} residual correlation; \texttt{furrr}/\texttt{purrr} provide the tidyverse collection; \texttt{data.table} provides the original-extended collection.
\end{itemize}

\textbf{Two-level seeding ensures reproducibility: driver-level seeds govern parameter grids; replicate-level seeds derive from a master seed via the replicate index.}

\begin{itemize}\setlength{\itemsep}{2pt}\setlength{\parskip}{0pt}
\item Driver level: \texttt{set.seed(42 + sweep\_id)} before parameter grid construction.
\item Replicate level: \texttt{set.seed(20260507)} at the loop top, with \texttt{parallel::mc.set.seed} inside \texttt{mclapply}.
\item The two seeding stages operate at distinct pipeline points and do not collide.
\end{itemize}

\textbf{All file paths for reproducing this manuscript are fixed: per-replicate RDS, ADEMP summary, driver scripts, pre-registered plan, and manuscript source.}

\begin{itemize}\setlength{\itemsep}{2pt}\setlength{\parskip}{0pt}
\item Per-replicate outputs: \texttt{analysis/data/quick-sim/05-design-replicates.rds}; ADEMP summary: \texttt{05-design-summary.txt}.
\item Driver scripts: \texttt{analysis/scripts/nof1-design-sensitivity/} and \texttt{analysis/scripts/quick-sim/05-design-prototype.R}.
\item Pre-registered ADEMP plan: \texttt{00-ademp-pre-registration.md}; manuscript source: \texttt{analysis/report/05-nof1-design-sensitivity/}.
\end{itemize}

\end{document}
